\section{Background and Related Work}
\label{sec:background}

\subsection{Obfuscation and virtualization}

Software obfuscation is the systematic transformation of a program to make its
behavior harder to recover while preserving semantics. Collberg, Thomborson, and
Low classified transformations into layout, data, control, and preventive
families \cite{collberg1997,collberg2002}. Control-flow flattening replaces a
structured CFG with a dispatcher over state variables; data transformations
rewrite expressions and constants; preventive transformations detect or impede
debuggers and analysis tools. Banescu and Pretschner give a modern tutorial and
taxonomy \cite{banescu2018}.

Code virtualization is a control transformation in which the protected function
is translated into a bytecode program for a private interpreter. The interpreter
and bytecode are embedded in the binary. The original native instructions can be
removed, and the bytecode can be encrypted, shuffled, and diversified. The
approach is attractive because it changes the attacker's problem from reading
native instructions to recovering both the ISA and the program. It is expensive
because every virtual instruction requires dispatch, decode, and interpretation.

Barak \etal{} proved that general-purpose obfuscation with strong
virtual-black-box security is impossible for arbitrary programs
\cite{barak2012}. The result does not rule out practical protection: it rules out
a universal transformation that hides all properties of every program. Sirus
therefore targets cost asymmetry and build-specific diversity, not
indistinguishability.

\subsection{Commercial protectors}

VMProtect is a commercial protector with virtualization, mutation, packing,
anti-debug, and licensing features \cite{vmprotect}. Code Virtualizer and Themida
from Oreans provide selective virtualization, multiple VM architectures, packing,
and anti-debug controls \cite{codevirtualizer,themida}. Denuvo Anti-Tamper is a
game-oriented product with its own virtual machine and integrity checks
\cite{denuvo}. These products are mature and support much broader instruction
sets and platforms than Sirus. Their internals are not public, and their security
claims are vendor-documented rather than independently reproduced in this work.

The comparison in this paper is deliberately narrow. We use VMProtect 3.11 Linux
demo as a measured performance peer because it can virtualize the same benchmark
function on the same host. We do not claim equivalence of security profiles,
packing, anti-debug, licensing, or platform coverage.

\subsection{Research systems}

Tigress is a source-to-source C transformer with a broad set of diversification
transformations, including virtualization, superoperators, reentrant interpreters,
anti-taint, and anti-aliasing passes \cite{tigress}. It is the closest research
peer in ambition, but it operates at source level and is not a binary protector.
Obfuscator-LLVM implements control-flow flattening, instruction substitution, and
bogus control flow at the LLVM IR level \cite{junod2015ollvm}; it is not a
virtualizer. MOVfuscator compiles programs into `mov'-only code as a research
demonstration \cite{domas2015movfuscator}; its cost profile is intentionally
extreme and it is not a selective production protector.

Sirus differs from these systems in three ways: it operates directly on ELF and
PE binaries, it is implemented in safe Rust with no third-party runtime crates,
and it ships an evidence harness that executes the protected binaries and the
external analysis tools. The trade-off is a strict instruction subset and a
narrow platform scope.

\subsection{Mixed Boolean-arithmetic obfuscation}

Mixed Boolean-arithmetic (MBA) expressions combine bitwise operations with
integer addition, subtraction, and multiplication. Zhou \etal{} showed that such
expressions can hide simple arithmetic and proposed rewrite-based
simplification \cite{zhou2007}. Modern simplifiers exploit the fact that a linear
MBA expression over $n$ variables is determined by its values on the $2^n$
Boolean inputs, and use linear algebra over the bit-vector space to recover a
compact form. GAMBA generalizes the approach to selected non-linear expressions
\cite{reichenwallner2023gamba}; SiMBA specializes in linear expressions
\cite{reichenwallner2022simba}; MBA-Blast is another widely used
simplifier \cite{mba_blast}.

Sirus uses a layered approach. Exact linear identities provide a baseline
rewrite, and one non-linear wrapper per function adds variable shifts and a
data-dependent modular inverse. The external simplifier study in
Section~\ref{sec:evaluation} reports both the reductions and the failures. The
result is measured cost asymmetry, not a claim that the expressions are
irreducible.

\subsection{Binary analysis and symbolic execution}

Ghidra \cite{ghidra} and radare2 \cite{radare2} are open reverse-engineering
frameworks used for static decompilation and disassembly. angr is a Python
binary-analysis platform with symbolic execution, constraint solving, and data
flow analysis \cite{shoshitaishvili2016angr}. KLEE is a symbolic execution engine
for LLVM bitcode that generates high-coverage test inputs and checks memory
safety \cite{cadar2008klee}. Z3 is the SMT solver commonly used by these tools
\cite{z3_2008}. PIN and DynamoRIO are dynamic instrumentation frameworks that can
trace a running process at instruction granularity \cite{pin2005}.

Sirus does not attempt to defeat these tools in the sense of a formal
impossibility result. It increases the size of the state space through encrypted
bytecode, randomized state identifiers, flattened control flow, MBA, and a
bounded fuel budget. The evaluation records what the tools did within fixed
budgets and treats a timeout as bounded resistance.

\subsection{Constant-time execution and timing measurement}

Timing attacks exploit data-dependent execution time, memory access patterns, or
microarchitectural state \cite{kocher1996,bernstein2005}. Spectre-class attacks
exploit speculative execution and are not solved by branchless source code
alone \cite{kocher2019spectre}. Constant-time programming therefore targets a
specific class of leakage: the sequence of instructions and memory addresses,
not every microarchitectural effect.

dudect-bencher is a statistical tool that measures whether two input classes are
distinguishable by execution time \cite{reparaz2017dudect}. Cachegrind, part of
Valgrind, provides deterministic instruction and cache-access counts
\cite{nethercote2007valgrind}. Sirus uses both: dudect for wall-clock
distinguishability and Cachegrind for input-independent instruction and address
counts on a controlled leaf function.

\subsection{Rust safety and verification}

Rust's ownership and borrowing rules provide memory and thread safety without a
garbage collector; RustBelt gives a formal foundation for the safe
fragment \cite{jung2018rustbelt,rustref}. Sirus additionally enforces
\code{forbid(unsafe_code)} in the safe tier, forbids heap allocation in the core
and VM, and uses fixed-capacity data structures. Kani provides bounded model
checking for Rust \cite{kani}, Miri detects undefined behavior in the Rust
abstract machine \cite{miri}, and Creusot/Why3 supports deductive verification of
Rust functions \cite{creusot2022,why3_2013}. These tools cover different scopes;
the evaluation states each scope explicitly.

\subsection{Positioning}

Table~\ref{tab:related} summarizes the qualitative position. Commercial entries
are based on public vendor documentation and are marked accordingly; they are not
independently verified in this work.

\begin{table}[htbp]
\centering
\small
\caption{Qualitative positioning. ``Y'' denotes documented or implemented,
``P'' partial or subset, and ``N'' not supported. Commercial entries are
vendor-documented capabilities, not independent measurements.}
\label{tab:related}
\begin{tabularx}{\textwidth}{@{}lcccccc@{}}
\toprule
Capability & Sirus & VMProtect & Code Virtualizer & Themida & Tigress & OLLVM \\
\midrule
Selective virtualization & Y & Y & Y & Y & Y & N \\
Full x86-64 ISA & N & Y & Y & Y & via C & via IR \\
Linux ELF runtime & Y & Y & Y & N & via C & Y \\
Windows PE runtime & P & Y & Y & Y & via C & via C \\
Multiple VM designs & N & Y & Y & Y & Y & N \\
Packing & N & Y & N & Y & N & N \\
Anti-debug / anti-VM & P & Y & P & Y & research & N \\
Runtime integrity & P & Y & P & Y & research & P \\
Safe Rust runtime & Y & N & N & N & N & N \\
Zero runtime crates & Y & N & N & N & N & N \\
Source available & Y & N & N & N & limited & Y \\
Local evidence harness & Y & N/A & N/A & N/A & papers & public \\
\bottomrule
\end{tabularx}
\end{table}

Sirus is best described as a reproducible, selective Linux-first virtualizer
with a Windows leaf-function line, within roughly one order of magnitude of a
commercial per-call VM on the same benchmark, but without the platform coverage,
anti-debug, packing, licensing, or maturity of commercial products. It is not a
replacement for a commercial protector on arbitrary Windows software.
